%------------------------------------------------------------------------------%

\usepackage{integracao_e_series}

\title{Séries Numéricas -- Propriedades}

%------------------------------------------------------------------------------%
\begin{document}

\addtitle

%------------------------------------------------------------------------------%
\section{Propriedades}

%------------------------------------------------------------------------------%
\begin{frame}{Linearidade das Séries}
  Se \(\sum a_n\) e \(\sum b_n\)
  são \emph{séries convergentes} e $c\in\R$,
  \pause
  então

  \vspace{\baselineskip}
  \begin{enumerate}[<+->]
    \setlength\itemsep{1em}
    \item \( \sum_{n=1}^\infty (a_n+b_n) =  \sum_{n=1}^\infty a_n +\sum_{n=1}^\infty b_n \)
    \item \( \sum_{n=1}^\infty (a_n-b_n) =  \sum_{n=1}^\infty a_n -\sum_{n=1}^\infty b_n \)
    \item \( \sum_{n=1}^\infty ka_n      = k\sum_{n=1}^\infty a_n \)
  \end{enumerate}
\end{frame}

%------------------------------------------------------------------------------%
\begin{frame}{Ordem dos termos}
  Não podemos alterar a ordem em que os termos são somados

  \pause
  \[
    1-1+1-1+1-1+1-1+1-1+\cdots \text{ diverge}
  \]

  \pause
  \[
    (1-1)+(1-1)+(1-1)+(1-1)+(1-1)+\cdots = 0
  \]

  \pause
  \[
    1-(1+1)-(1+1)-(1+1)-(1+1)-(1+1)-\cdots = -\infty
  \]
\end{frame}

%------------------------------------------------------------------------------%
\begin{frame}{Verificando a Convergência}
  Basta analisar a soma dos termos a partir de um índice $N$

  \pause
  \vspace{0.5\baselineskip}
  Se
  \[
    \sum_{n={\color{blue}N}}^\infty a_n
  \]
  converge, ou diverge, então
  \[
    \sum_{n={\color{blue}0}}^\infty a_n
  \]
  também
\end{frame}

%------------------------------------------------------------------------------%
\begin{frame}{Alterando os Termos de uma Série}
  Adicionar, retirar ou permutar um \emph{número finito de termos} em uma série,\\
  não altera o fato da série convergir ou divergir

  \pause
  \vspace{\baselineskip}
  Porém, pode alterar o valor da soma da série, caso ele exista
\end{frame}

%------------------------------------------------------------------------------%
\begin{frame}{\contentsname}
  \tableofcontents
\end{frame}

%------------------------------------------------------------------------------%
\end{document}
%------------------------------------------------------------------------------%
